\section{Discriminant evolution for a finite cluster of arbitrary size}
\label{m:section}

The two-root equation has a regular extension to every finite cluster,
but its constant forcing does not.  This distinction gives both a useful
general identity and exact limits on what can be inferred from two real
cap signs.  The results in this section develop the analytic statements
recorded in \cite{note09}; their earlier internal review is \cite{review08}.

\subsection{A regular equation through collisions}

Let $D$ be a bounded Jordan complex-time domain with piecewise smooth
boundary, and let $H$
be holomorphic near
$\overline D\times\{z:|z-a|\leq R\}$, with $H_t=-H_{zz}$.
Assume that, for every $t\in\overline D$, $H(t,\cdot)$ has exactly
$m\geq2$ zeros in the outer disk, all in $|z-a|<r<R$, and has no zero
on the outer circle.  These assumptions are understood with strict
margins sufficient to extend them to a neighborhood of the closed
product.  The cluster factorization gives
\begin{equation}\label{m:factor}
 P(t,z)=\prod_{i=1}^m(z-z_i(t)),\qquad H(t,z)=U(t,z)P(t,z),
 \qquad q=\frac{U_z}{U},\qquad f=\frac{H_z}{H},
\end{equation}
where $P$ has single-valued holomorphic coefficients and $U$ is
holomorphic and nonvanishing on the spatial disk.  Individual root
labels need not be single-valued.  In particular, $q$ includes the
complete exterior contribution, including any zero-free factor of $H$.

Put $N=m(m-1)/2$ and define
\begin{equation}\label{m:DT}
 \mathfrak D=\prod_{i<j}(z_i-z_j)^2,\qquad
 \mathfrak T=\sum_{i<j}
       \prod_{\substack{k<\ell\\(k,\ell)\ne(i,j)}}
                         (z_k-z_\ell)^2.
\end{equation}
The empty product is one.  Both expressions are symmetric polynomials
in the roots and hence are holomorphic functions of the coefficients
of $P$.  Thus \eqref{m:DT} defines $\mathfrak T$ at collisions; the
quotient expression
$\mathfrak T=\mathfrak D\sum_{i<j}(z_i-z_j)^{-2}$ is used only when
the roots are distinct.

\begin{theorem}[Regular finite-cluster discriminant equation]
\label{m:evolution}
Under these assumptions,
\begin{align}
 \mathfrak D'&=8\mathfrak T+B\mathfrak D,
                                                        \label{m:ode}\\
 B(t)&=\frac{2}{2\pi i}\oint_{|\zeta-a|=R}
        \frac{H_z(t,\zeta)}{H(t,\zeta)}
        \frac{P_{zz}(t,\zeta)}{P(t,\zeta)}\,d\zeta.
                                                        \label{m:B}
\end{align}
The coefficient $B$ is holomorphic through every internal
collision.  For any holomorphic scalar $\kappa(t)$, replacing $H_z$
by $H_z-\kappa H$ in \eqref{m:B} leaves it unchanged.  If, on the
complete outer wall,
\begin{equation}\label{m:wallinputs}
 |H|\geq\delta>0,\qquad |H_z-\kappa H|\leq L_\kappa,
\end{equation}
then
\begin{equation}\label{m:Bbound}
 |B|\leq
 \frac{2R\,m(m-1)L_\kappa}{\delta(R-r)^2}.
\end{equation}
In particular, the bound contains no inverse internal root separation.
\end{theorem}

\begin{proof}
At a simple zero, differentiation and the heat equation give
\begin{equation}\label{m:motion}
 z_i'=\frac{H_{zz}(t,z_i)}{H_z(t,z_i)}
      =2\sum_{j\ne i}\frac{1}{z_i-z_j}+2q(t,z_i).
\end{equation}
Writing $d_{ij}=z_i-z_j$, we obtain
\begin{equation}\label{m:pairmotion}
 \frac{z_i'-z_j'}{d_{ij}}
 =\frac{4}{d_{ij}^2}
 -2\sum_{k\ne i,j}\frac{1}{(z_i-z_k)(z_j-z_k)}
 +2\frac{q(t,z_i)-q(t,z_j)}{d_{ij}}.
\end{equation}
The middle terms cancel after summation over all pairs, because every
unordered triple contributes a multiple of the identity
\[
 \frac{1}{(x-z)(y-z)}+
 \frac{1}{(x-y)(z-y)}+
 \frac{1}{(y-x)(z-x)}=0.
\]
Logarithmic differentiation therefore yields
\begin{equation}\label{m:logderivative}
 \frac{\mathfrak D'}{\mathfrak D}
 =8\sum_{i<j}\frac{1}{d_{ij}^2}
  +4\sum_{i<j}\frac{q(t,z_i)-q(t,z_j)}{z_i-z_j}.
\end{equation}
The divided differences in the last term extend holomorphically when
their arguments coincide.  Cauchy's formula and
\[
 \frac{P_{zz}(t,\zeta)}{P(t,\zeta)}
    =2\sum_{i<j}\frac{1}{(\zeta-z_i)(\zeta-z_j)}
\]
give
\[
 4\sum_{i<j}\frac{q(t,z_i)-q(t,z_j)}{z_i-z_j}
 =\frac{2}{2\pi i}\oint_{|\zeta-a|=R}
             q(t,\zeta)\frac{P_{zz}(t,\zeta)}{P(t,\zeta)}\,d\zeta.
\]
Replacing $q$ by $f-P_z/P$ introduces the rational integral of
$P_zP_{zz}/P^2$.  All finite poles of this rational function lie inside
the contour, and it is $O(\zeta^{-3})$ at infinity.  Its total residue
is consequently zero.  This proves \eqref{m:B}.  The same reasoning
applied to $P_{zz}/P=O(\zeta^{-2})$ proves the exact scalar subtraction.

Every denominator in \eqref{m:B} is nonzero on the integration circle,
so the coefficient is holomorphic in time.  The simple-root set is
dense: at any time, the local heat splitting of a zero of multiplicity
$k$ has distinct Hermite leading roots after a nonzero small time
increment.  The finitely many clusters in that fiber can all be
separated by one sufficiently small increment.  Thus $\mathfrak D$
is not identically zero.  Multiplying \eqref{m:logderivative} by
$\mathfrak D$ proves \eqref{m:ode} on a dense set, and then everywhere
by holomorphic continuation.

Finally, $|\zeta-z_i|\geq R-r$ on the circle, whence
\[
 \left|\frac{P_{zz}}{P}\right|
       \leq\frac{m(m-1)}{(R-r)^2}.
\]
Using the contour length $2\pi R$ and \eqref{m:wallinputs} in
\eqref{m:B} proves \eqref{m:Bbound}.
\end{proof}

For $m=2$ one has $\mathfrak T=1$, so the theorem is exactly the
two-root evolution equation.  For larger clusters, it keeps every
root in the birth while retaining the same power two of the spatial
collar in the coefficient bound.

\subsection{An explicit window at a hypothesized full birth}

For the complete theta source, write
\[
 B_j=M_j(T_0,R_0)
   =\int_0^\infty\Phi(u)u^j e^{T_0u^2}\cosh(R_0u)\,du,
 \qquad T_0,R_0>0.
\]
These finite positive moments bound all derivatives needed below.
\begin{proposition}[Parameterized geometry at multiplicity $m$]
\label{m:window}
Suppose $a\in\mathbb R$ and $H(0,\cdot)$ has exact multiplicity
$m\geq2$ at $a$.  Set
\begin{align}
 q_m&=\frac{|\partial_z^mH(0,a)|}{m!}>0,
 &R&=\min\left\{\frac{R_0}{2},
                \frac{(m+1)!q_m}{4B_{m+1}}\right\},
 &r&=\frac R2,                                             \label{m:radii}\\
 h&=\min\left\{\frac{T_0}{2},
                \frac{q_mr^m}{4\sqrt2 B_2}\right\}.
                                                               \label{m:width}
\end{align}
For $|\operatorname{Re}t|,|\operatorname{Im}t|\leq h$, each of the
two disks of radii $r$ and $R$ contains exactly $m$ roots, and
\begin{equation}\label{m:paidwall}
 |H(t,z)|\geq\frac{q_ms^m}{2}
       \quad\hbox{if } |z-a|=s,\quad s\in\{r,R\}.
\end{equation}
The complete coefficient in \eqref{m:B} obeys
\begin{equation}\label{m:paidB}
 |B|\leq
       \frac{16m(m-1)B_1}{q_mR^{m+1}}.
\end{equation}
\end{proposition}

\begin{proof}
Put $c_m=\partial_z^mH(0,a)/m!$.  Spatial Taylor expansion, with the
complete derivative bound on the segment from $a$ to $z$, gives
\[
 |H(0,z)-c_m(z-a)^m|
 \leq\frac{B_{m+1}s^{m+1}}{(m+1)!}
 \leq\frac{q_ms^m}{4} \qquad (|z-a|=s\in\{r,R\}).
\]
Along the straight time segment, $|H_t|=|H_{zz}|\leq B_2$ and
$|t|\leq\sqrt2h$, so
\[
 |H(t,z)-H(0,z)|\leq\sqrt2hB_2\leq\frac{q_mr^m}{4}.
\]
The total error is at most $q_ms^m/2$, strictly less than
$|c_m(z-a)^m|$.  Rouch\'e's theorem gives both root counts and the
triangle inequality gives \eqref{m:paidwall}.  The inequalities have
strict slack relative to the zero-avoidance and root-count requirements,
so the contract extends near the closed product.  In
\eqref{m:Bbound}, take $\kappa=0$, $L_\kappa=B_1$,
$\delta=q_mR^m/2$ and $R-r=R/2$, obtaining \eqref{m:paidB}.
\end{proof}

This construction is conditional on the event and its nonzero order-$m$
derivative.  It gives a positive window at every such finite event,
with all theta modes retained in $B_j$.  It supplies neither a sign of
the discriminant nor nonvanishing on the complex time boundary.
The dependence on $q_m$ and $m$ is explicit and may deteriorate with
height.

\subsection{Birth order, forcing, and charge}

Suppose the entire cluster coalesces at $(t_0,z_0)$.  Write
$s=t-t_0$ and $s=\rho^2$.  The local heat splitting is
\begin{equation}\label{m:hermitesplit}
 z_i(t_0+\rho^2)=z_0+\sqrt2\lambda_i\rho+O(\rho^2),
\end{equation}
where $\lambda_1,\ldots,\lambda_m$ are the distinct real roots of the
monic probabilists' Hermite polynomial $\operatorname{He}_m$.
Indeed, the leading part of parabolic weight $m$ is
\[
 e^{-s\partial_z^2}(z-z_0)^m
    =(2s)^{m/2}\operatorname{He}_m
                   \left(\frac{z-z_0}{\sqrt{2s}}\right).
\]
All higher initial spatial Taylor terms have larger parabolic weight;
analytic continuation of the simple scaled roots gives
\eqref{m:hermitesplit}.

\begin{corollary}[Order at a full birth]\label{m:birth}
At a full multiplicity-$m$ birth,
\begin{align}
 \mathfrak D(t)&=d_ms^N+O(s^{N+1}),
 &d_m&=2^N\prod_{i<j}(\lambda_i-\lambda_j)^2>0,
                                                               \label{m:Dleading}\\
 \mathfrak T(t)&=\frac{Nd_m}{8}s^{N-1}+O(s^N).
                                                               \label{m:Tleading}
\end{align}
If $\mathfrak D$ is nonzero on $\partial D$, the product-domain
charge is
\begin{equation}\label{m:charge}
 Q=\frac{1}{2\pi i}\oint_{\partial D}
                   \frac{\mathfrak D'}{\mathfrak D}\,dt
  =\frac{8}{2\pi i}\oint_{\partial D}
                   \frac{\mathfrak T}{\mathfrak D}\,dt.
\end{equation}
\end{corollary}
\begin{proof}
Taking the product of squared differences in \eqref{m:hermitesplit}
gives the leading coefficient in \eqref{m:Dleading}.  Holomorphy in
$s$ removes any noninteger power from its remainder.  Substitution
in \eqref{m:ode}, using holomorphy of $B$, proves
\eqref{m:Tleading}.  The first equality in \eqref{m:charge} is the
finite-cluster charge formula.  For the second, integrate
\eqref{m:ode} divided by $\mathfrak D$; the integral of the holomorphic
coefficient $B$ is zero.
\end{proof}

The exterior contribution still changes $\mathfrak D$, $\mathfrak T$
and their zero locations through \eqref{m:ode}; its vanishing contour
integral is not permission to omit it from the evolution.  For $m>2$
the forcing vanishes at a full birth.  This removes the constant
forcing used by the two-root univalence argument.  A full birth has
charge $N$, while two real discriminant signs record at most its
parity.  Conjugate nonreal event pairs can contribute further even
charge without changing those signs.

\subsection{A cubic real-time monotonicity theorem}

Assume now that $a$ is real, $D$ is invariant under conjugation and
$H(\overline t,\overline z)=\overline{H(t,z)}$.  For $m=3$, depress
the cluster polynomial:
\[
 P(t,z)=(z-c(t))^3+p(t)(z-c(t))+q_0(t).
\]
The notation $q_0$ distinguishes the constant coefficient from $U_z/U$.
Direct expansion gives
\begin{equation}\label{m:cubicDT}
 \mathfrak D=-4p^3-27q_0^2,\qquad \mathfrak T=9p^2.
\end{equation}
For the second identity, $\mathfrak T$ is a translation-invariant
symmetric polynomial of degree four.  After centering, the fundamental
symmetric polynomials have degrees two and three, so it is a constant
multiple of $p^2$.  Substituting roots $(-1,0,1)$ gives the constant
$9$.

\begin{proposition}[At most one real cubic collision time]
\label{m:cubicmonotone}
On each real interval contained in $D$, a proper real-symmetric
three-root heat cluster has at most one collision time.  An
independently located collision at a strictly negative time in such an
interval excludes an additional collision at time zero in that
interval.
\end{proposition}
\begin{proof}
Choose an antiderivative of $B$ which is real on the real
interval, and put $\mu=\exp(-\int B)$.  Then $\mu>0$ there
and \eqref{m:ode}--\eqref{m:cubicDT} imply
\begin{equation}\label{m:cubicderivative}
                  (\mu\mathfrak D)'=72\mu p^2\geq0.
\end{equation}
If $p$ is not identically zero, its real zeros are isolated.  The
integral of the right side over every nontrivial subinterval is then
positive, so $\mu\mathfrak D$ is strictly increasing and has at most
one zero.  If $p$ is identically zero, $\mu\mathfrak D$ is constant
and nonzero, since $\mathfrak D$ cannot be identically zero.  There
are no collisions in that case.
\end{proof}

A real event in this cubic cluster can be a double, of charge one,
or a triple, of charge three.  Thus the proposition bounds the number
of real event locations, not their total charge.  It also gives no
bound on nonreal event times.  At complex time $p^2$ need not have
positive real part, even if $\mu$ lies in a fixed open half-plane.
For $m\geq4$, even the real inequality in
\eqref{m:cubicderivative} need not persist.

\subsection{Exact controls for the failure of the cap formula}

The following polynomial heat solutions pay the cluster and boundary
conditions explicitly.  They disprove the extension of the two-root
cap formula to arbitrary multiplicity, without challenging its
two-root hypothesis.  Their conclusions are algebraic; the exact
arithmetic replay in \cite{controls} is supplementary support.

\begin{proposition}[A cubic with three event times]
\label{m:cubiccontrol}
Let
\[
 H(t,z)=z^3-6tz+2,\qquad
 D=\{t:|\operatorname{Re}t|<1,\ |\operatorname{Im}t|<1\},
 \qquad r=4,\quad R=5,\quad a=0.
\]
This is a proper three-root heat cluster with $B=0$.  Its
total charge is three, although
$\bigl(\operatorname{sgn}\mathfrak D(1)
       -\operatorname{sgn}\mathfrak D(-1)\bigr)/2=1$.
\end{proposition}
\begin{proof}
One has $H_t=-6z=-H_{zz}$.  On the inner wall and closed time
rectangle,
\[
 |6tz-2|\leq24\sqrt2+2<64=|z^3|.
\]
Rouch\'e's theorem puts all three polynomial roots strictly in the
inner disk, hence also in the outer disk; both walls are nonzero.
Here $U=1$, and thus $B=0$.  Formula
\eqref{m:cubicDT} gives
\[
 \mathfrak D(t)=864t^3-108,
 \qquad \mathfrak T(t)=324t^2.
\]
The event times are the three simple roots of $t^3=1/8$:
\[
 \frac12,\qquad -\frac14+\frac{\sqrt3}{4}i,
 \qquad -\frac14-\frac{\sqrt3}{4}i.
\]
At each event the repeated root is $z=1/(2t)$, of modulus one;
$H_{zz}=6z\ne0$ there, so each event is an ordinary double.  All
three lie strictly within the product domain, and no event is on its
boundary.  Each has charge one.  Finally,
$\mathfrak D(-1)=-972$ and $\mathfrak D(1)=756$, proving the
claimed discrepancy.
\end{proof}

\begin{proposition}[A quadruple endpoint with equal positive cap signs]
\label{m:quarticcontrol}
On the same time rectangle, let
\[
 H(t,z)=z^4-12tz^2+12t^2,
 \qquad a=0,\quad r=5,\quad R=6.
\]
This proper four-root heat cluster has $B=0$, nonzero
time-boundary discriminant, equal positive real cap discriminants and
total charge six.  For every positive real time its spatial roots are
real and simple.
\end{proposition}
\begin{proof}
The heat equation follows by direct differentiation.  The inner-wall
estimate is
\[
 |12tz^2-12t^2|\leq300\sqrt2+24<625=|z^4|,
\]
so all four roots lie strictly in the inner disk throughout the
closed rectangle.  Again $U=1$ and $B=0$.  Set
$\alpha=6+2\sqrt6$, $\beta=6-2\sqrt6$.  The factorization
\[
 H(t,z)=(z^2-\alpha t)(z^2-\beta t)
\]
gives roots $\pm\sqrt{\alpha t}$ and $\pm\sqrt{\beta t}$.
The squared differences within the two opposite pairs contribute
$16\alpha\beta t^2$, and the four cross-pair squared differences
contribute $(\alpha-\beta)^4t^4$.  Therefore
\[
 \mathfrak D(t)=16\alpha\beta(\alpha-\beta)^4t^6
              =1{,}769{,}472\,t^6,
 \qquad
 \mathfrak T(t)=1{,}327{,}104\,t^5,
\]
where the last identity also follows from $\mathfrak D'=8\mathfrak T$.
The only event is $(0,0)$, of charge six.  Both cap discriminants are
positive, and no point of the time boundary is an event.  Since
$\alpha>\beta>0$, the four displayed roots are real and distinct for
$t>0$.  Notice also that $\mathfrak T$ changes sign on real time.
\end{proof}

More generally, the Hermite family $e^{-t\partial_z^2}z^m$ has
$\mathfrak D=d_mt^N$, $\mathfrak T=Nd_mt^{N-1}/8$ and
$B=0$; a sufficiently large spatial disk gives a proper
finite window.  Thus zero exterior phase alone is compatible with
every full-birth charge.  The controls above concern the specified
phase-and-cap inference.  They do not preclude other boundary
estimates or estimates using additional properties of the theta
source.

\subsection{Why centered variance cannot replace the discriminant}

Define
\[
 c=\frac1m\sum_i z_i,\qquad
 V=\sum_i(z_i-c)^2=\frac1m\sum_{i<j}(z_i-z_j)^2.
\]
The equality follows by expanding the pairwise sum.  This quantity is
holomorphic even when root labels are not.

\begin{proposition}[Centered second-moment evolution]\label{m:variance}
For a proper $m$-root heat cluster,
\begin{equation}\label{m:Vflow}
 V'=2m(m-1)+4\sum_i(z_i-c)q(t,z_i).
\end{equation}
At a full birth $(t_0,z_0)$,
\begin{equation}\label{m:Vbirth}
 V=2m(m-1)(t-t_0)+O((t-t_0)^2).
\end{equation}
\end{proposition}
\begin{proof}
Since $\sum_i(z_i-c)=0$, differentiating gives
$V'=2\sum_i(z_i-c)z_i'$.  Substitute \eqref{m:motion}.  For
each unordered internal pair the contribution is
\[
 4\left(\frac{z_i-c}{z_i-z_j}
       +\frac{z_j-c}{z_j-z_i}\right)=4.
\]
There are $N$ pairs, proving \eqref{m:Vflow} at simple fibers.  The
remaining symmetric sum is holomorphic, for it equals
\[
 \frac{1}{2\pi i}\oint_{|\zeta-a|=R}
       (\zeta-c)q(t,\zeta)\frac{P_z(t,\zeta)}{P(t,\zeta)}\,d\zeta.
\]
Continuation proves the identity everywhere.  At a full birth $V=0$
and each $z_i-c=0$, so evaluating \eqref{m:Vflow} gives
\eqref{m:Vbirth}.
\end{proof}

The real sign transition in \eqref{m:Vbirth} is compatible with
every Hermite birth.  Moreover, $V$ need not have the discriminant's
zero set: the cubic in Proposition~\ref{m:cubiccontrol} has
$V=12t$, hence $V(0)=0$ but $\mathfrak D(0)=-108\ne0$.
On an all-real fiber, $V=0$ does force full coalescence, but that
fiber condition cannot be assumed on an unknown negative cap.
Conversely, $V>0$ does not exclude a partial collision in a larger
real cluster.  Neither centered variance nor two discriminant signs
therefore replace the full multiplicity-sensitive event calculation.
